\documentclass[11pt,a4paper]{book}

\usepackage[T1]{fontenc}
\usepackage{lmodern}
\usepackage{amsmath,amssymb,amsthm,mathtools,bm}
\usepackage[margin=2.6cm]{geometry}
\usepackage{enumitem}
\usepackage{array}
\usepackage[colorlinks=true,linkcolor=black,citecolor=black,urlcolor=black]{hyperref}

\numberwithin{equation}{chapter}
\setlist{itemsep=2pt,topsep=4pt}

% --- macros -----------------------------------------------------------
\newcommand{\dd}{\mathrm{d}}
\newcommand{\ii}{\mathrm{i}}
\newcommand{\ee}{\mathrm{e}}
\newcommand{\vk}{\bm{k}}
\newcommand{\vr}{\bm{r}}
\newcommand{\vp}{\bm{p}}
\newcommand{\Fhyp}{{}_1F_1}
\newcommand{\sgn}{\operatorname{sgn}}
\newcommand{\artanh}{\operatorname{artanh}}
\newcommand{\Imag}{\operatorname{Im}}
\newcommand{\Real}{\operatorname{Re}}
\newcommand{\dsdO}{\frac{\dd\sigma}{\dd\Omega}}
\newcommand{\braket}[2]{\langle #1|#2\rangle}
\newcommand{\bra}[1]{\langle #1|}
\newcommand{\ket}[1]{|#1\rangle}
\newcommand{\afs}{\alpha_{\mathrm{f}}}

\title{Scattering by Long-Range Forces:\\[2pt]
       Classical, Quantum and Relativistic Theory}
\author{Draft chapters for a physics text}
\date{}

\begin{document}
\frontmatter
\maketitle
\tableofcontents

\chapter*{Scope and conventions}
\addcontentsline{toc}{chapter}{Scope and conventions}

Most treatments of scattering rest on one of two simplifications. Either the
interaction vanishes beyond a finite radius, so that partial waves, phase
shifts, scattering lengths, resonances and hard spheres apply; or the
interaction is weak enough to be expanded in its strength, giving the Born
series, Fermi's golden rule and Feynman diagrams. These chapters omit both on
purpose. What remains is narrower in scope but richer in structure:
\begin{itemize}
  \item the language of scattering that does not depend on the dynamics:
        cross sections, reference frames, invariants;
  \item exact treatments of infinite-range forces, above all the Coulomb
        potential, where finite-range methods fail and the classical, quantum
        and relativistic descriptions can be compared in closed form;
  \item non-perturbative structural results: unitarity, identical particles,
        the meaning of asymptotic states, analyticity.
\end{itemize}

Part~I (Chapters~1--2) sets up kinematics and the classical theory.
Part~II (Chapters~3--4) treats non-relativistic quantum scattering and ends
with the exact Coulomb amplitude. Part~III (Chapters~5--8) treats relativistic
scattering: classical orbits, the Klein--Gordon and Dirac Coulomb problems, the
Klein step, the invariant formulation, and a final comparison of the three
viewpoints.

\paragraph{Conventions.}
\begin{itemize}
  \item The Coulomb strength is $\alpha=Z_1Z_2e^2$ in Gaussian units (in SI
        replace $e^2$ by $e^2/4\pi\varepsilon_0$). The potential is
        $V(r)=\alpha/r$, so $\alpha>0$ is repulsive and $\alpha<0$ attractive.
  \item $E$ is the kinetic energy in the centre-of-mass (CM) frame unless
        stated otherwise; $\mu$ is the reduced mass, $v$ the relative speed at
        infinity, $p=\mu v$ and $k=p/\hbar$.
  \item $\eta=\alpha/(\hbar v)$ is the (signed) Sommerfeld parameter.
        Relativistically $v=pc^2/E$.
  \item $\zeta=\alpha/\hbar c$ (equal to $Z\afs$ for a nucleus of charge $Z$)
        measures how relativistic the Coulomb problem is; note
        $\zeta=\eta\,v/c$.
  \item An asterisk marks CM quantities; the subscript L marks the laboratory
        frame (target at rest).
\end{itemize}

\mainmatter

% =====================================================================
\part{Foundations}
% =====================================================================

\chapter{The Scattering Experiment and Its Kinematics}

\section{Cross sections}

A scattering experiment sends a beam of projectiles, with flux $J$ (particles
per unit area per unit time, measured transverse to the beam), onto $N$ targets
and counts the particles emerging per unit time into the solid angle
$\dd\Omega$ around the direction $(\theta,\varphi)$. If the targets do not
shadow one another,
\begin{equation}\label{eq:flux}
  \dd\dot N = J\,N\,\dsdO\,\dd\Omega .
\end{equation}
The differential cross section $\dd\sigma/\dd\Omega$ has the dimension of an
area; the total cross section is $\sigma=\int(\dd\sigma/\dd\Omega)\,\dd\Omega$.
Three remarks fix the meaning of these definitions.
\begin{enumerate}
  \item A cross section is a ratio of rates. It is therefore independent of
        beam intensity, target thickness and detector efficiency, and encodes
        only the dynamics of the interaction.
  \item Classically, $\dd\sigma$ is literally an area: the part of the
        incident beam's transverse cross section whose particles end up in
        $\dd\Omega$.
  \item Because $\sigma$ is a transverse area, it is unchanged by Lorentz
        boosts along the beam axis. Total cross sections therefore agree
        between the lab and CM frames; differential cross sections differ only
        through the change of angle.
\end{enumerate}
The definition \eqref{eq:flux} assumes a beam that is uniform over a region much
wider than the range of the force. For an infinite-range force this assumption
can never be fully met, and we will see that it is the origin of the infinite
total cross sections that haunt Coulomb scattering.

\section{Two-body reduction and the lab frame (non-relativistic)}

For two particles interacting through a potential depending on
$\vr=\vr_1-\vr_2$, the motion separates into free motion of the centre of mass
and motion of a single particle of reduced mass $\mu=m_1m_2/(m_1+m_2)$ in the
potential $V(r)$. Dynamics determines $\dd\sigma/\dd\Omega^*$ in the CM frame;
everything else is kinematics.

Let a projectile $m_1$ hit a target $m_2$ at rest and put $\rho=m_1/m_2$. For
elastic scattering through CM angle $\theta^*$,
\begin{align}
  \tan\theta_{\mathrm L} &= \frac{\sin\theta^*}{\rho+\cos\theta^*}, \label{eq:labtan}\\[2pt]
  \left.\dsdO\right|_{\mathrm L} &=
      \frac{(1+\rho^2+2\rho\cos\theta^*)^{3/2}}{|1+\rho\cos\theta^*|}
      \left.\dsdO\right|^{*}, \label{eq:labxs}\\[2pt]
  T_2 &= \frac{4m_1m_2}{(m_1+m_2)^2}\,E_{\mathrm L}\sin^2\frac{\theta^*}{2}.
      \label{eq:recoil}
\end{align}
Here $E_{\mathrm L}$ is the projectile's lab kinetic energy and $T_2$ the recoil
energy of the target. For $\rho<1$ the lab angle covers
$0\le\theta_{\mathrm L}\le\pi$; for $\rho>1$ it cannot exceed
$\theta_{\max}$ with $\sin\theta_{\max}=1/\rho$, and each lab angle below
$\theta_{\max}$ corresponds to two CM angles. Equation~\eqref{eq:recoil} shows
that the energy transferred is a simple function of the CM angle alone, and it
will turn Rutherford's angular law into an energy-transfer law
(Section~\ref{sec:coullog}).

\section{Relativistic kinematics}

In this section $c=1$. For $p_1+p_2\to p_3+p_4$ define the Mandelstam
invariants
\begin{equation}\label{eq:mandelstam}
  s=(p_1+p_2)^2,\quad t=(p_1-p_3)^2,\quad u=(p_1-p_4)^2,\qquad
  s+t+u=m_1^2+m_2^2+m_3^2+m_4^2 .
\end{equation}
For elastic scattering ($m_3=m_1$, $m_4=m_2$) with the target at rest in the
lab:
\begin{itemize}
  \item $s=m_1^2+m_2^2+2m_2E_{\mathrm L}$, and the CM momentum is
        $p^*=m_2p_{\mathrm L}/\sqrt s$, where $p_{\mathrm L}$ is the beam
        momentum in the lab;
  \item the projectile's CM energy is $E_1^*=(s+m_1^2-m_2^2)/(2\sqrt s)$;
  \item $t=-2p^{*2}(1-\cos\theta^*)$, so $-4p^{*2}\le t\le 0$, and the target's
        lab recoil energy is $T_2=-t/(2m_2)$.
\end{itemize}
The CM moves in the lab with $\beta_c=p_{\mathrm L}/(E_{\mathrm L}+m_2)$ and
$\gamma_c=(E_{\mathrm L}+m_2)/\sqrt s$. The lab scattering angle follows from
transforming the projectile's momentum:
\begin{equation}\label{eq:lab-rel}
  \tan\theta_{\mathrm L}=\frac{\sin\theta^*}
        {\gamma_c\left(\cos\theta^*+\beta_c/\beta_1^*\right)},
  \qquad \beta_1^*=p^*/E_1^* .
\end{equation}
The non-relativistic limit is $\gamma_c\to1$, $\beta_c/\beta_1^*\to m_1/m_2$,
which reproduces \eqref{eq:labtan}. For equal masses $\beta_c=\beta_1^*$, so
$\tan\theta_{\mathrm L}=\tan(\theta^*/2)/\gamma_c$. The recoiling particle
emerges at $\tan\theta_{\mathrm L}'=\cot(\theta^*/2)/\gamma_c$, hence
\begin{equation}\label{eq:product}
  \tan\theta_{\mathrm L}\,\tan\theta_{\mathrm L}'=\frac1{\gamma_c^2}
  =\frac{2}{\gamma_{\mathrm L}+1},
\end{equation}
where $\gamma_{\mathrm L}$ is the projectile's lab Lorentz factor.
Non-relativistically the product is~1 and the two outgoing particles separate
at exactly $90^\circ$; relativistically the opening angle is always smaller
than $90^\circ$.

Two invariants will recur. The first is the flux factor,
$F=4\sqrt{(p_1\!\cdot\! p_2)^2-m_1^2m_2^2}=4p^*\sqrt s=4m_2p_{\mathrm L}$. The
second is the relation between the invariant $\dd\sigma/\dd t$ and the CM
angular cross section,
\begin{equation}\label{eq:dsdt}
  \frac{\dd\sigma}{\dd t}=\frac{\pi}{p^{*2}}\left.\dsdO\right|^{*}.
\end{equation}
Equation~\eqref{eq:dsdt} follows from $\dd t=2p^{*2}\,\dd(\cos\theta^*)$ and
$\dd\Omega=2\pi\,\dd(\cos\theta^*)$. It already hints at the main theme of
Chapter~7: the Rutherford law, $\dd\sigma/\dd\Omega\propto1/\sin^4(\theta^*/2)$,
is the statement $\dd\sigma/\dd t\propto1/t^2$.

\section{Problems}
\begin{enumerate}
  \item Derive \eqref{eq:labtan} and \eqref{eq:labxs} from the velocity-addition
        rule between the CM and lab frames.
  \item Prove \eqref{eq:product}, and show that for a head-on collision of
        equal masses the projectile stops.
  \item For the elastic scattering of a proton of lab kinetic energy 1~GeV on a
        proton at rest, find $s$, $p^*$, $\gamma_c$, and the largest lab angle
        at which both outgoing protons can be detected together.
\end{enumerate}

\section*{Notes and references}
Collision kinematics is developed concisely in \cite{LL1}. Relativistic kinematics, the Mandelstam variables and the K\"all\'en function are collected, with worked examples, in \cite{BK,Hagedorn} and in the kinematics review of \cite{PDG}; the invariants were introduced in \cite{Mandelstam58}. The definition of a cross section and its invariance under boosts are discussed in the general setting of scattering theory in \cite{Weinberg}.

% =====================================================================
\chapter{Classical Scattering}
% =====================================================================

\section{Orbits and the deflection function}

In the CM frame a central potential $V(r)$ makes the relative motion planar.
Let the incident relative speed be $v$ and the impact parameter $b$. Two
constants of motion fix everything: the energy $E=\tfrac12\mu v^2$ and the
angular momentum $L=\mu vb$. The radial and angular equations are
\[
  \dot r^2=\frac2\mu\Bigl[E-V(r)-\frac{L^2}{2\mu r^2}\Bigr],\qquad
  \dot\varphi=\frac{L}{\mu r^2}.
\]
Eliminating time, the angle swept between the two asymptotes is
\begin{equation}\label{eq:phi-swept}
  \Delta\varphi(b)=2\int_{r_0}^{\infty}
     \frac{b\,\dd r/r^2}{\sqrt{1-V(r)/E-b^2/r^2}},
\end{equation}
where $r_0$ is the largest root of the radicand (the distance of closest
approach). For $V=0$, $r_0=b$ and $\Delta\varphi=\pi$. We define the
deflection function
\begin{equation}\label{eq:deflection}
  \Theta(b)=\pi-\Delta\varphi(b),
\end{equation}
which is positive for repulsion and negative for attraction. The observed
scattering angle $\theta\in[0,\pi]$ is the value of $|\Theta+2\pi n|$ that falls
in that interval, for an integer $n$.

\section{From deflection function to cross section}

Particles with impact parameters between $b$ and $b+\dd b$ and azimuth between
$\varphi$ and $\varphi+\dd\varphi$ cross the beam's transverse plane at a rate
$J\,b\,\dd b\,\dd\varphi$. A central force conserves the azimuth, so those
particles leave into $\sin\theta\,\dd\theta\,\dd\varphi$. Equating the two, and
summing over all impact parameters $b_i$ that lead to the same $\theta$,
\begin{equation}\label{eq:xs-sum}
  \dsdO=\sum_i\frac{b_i}{\sin\theta}\left|\frac{\dd b_i}{\dd\theta}\right| .
\end{equation}
Three classical singularities of \eqref{eq:xs-sum} are generic and give the
cleanest bridge to quantum mechanics.
\begin{itemize}
  \item \textbf{Rainbow.} Where $\dd\Theta/\dd b=0$, $\dd b/\dd\theta$ is
        infinite and the classical cross section diverges (as
        $(\theta-\theta_r)^{-1/2}$). Two branches merge at a caustic; quantum
        mechanically the divergence becomes an Airy-function peak with
        supernumerary oscillations.
  \item \textbf{Glory.} Where $\Theta=0$ or $\pi$ at $b\neq0$, the factor
        $1/\sin\theta$ makes $\dd\sigma/\dd\Omega$ diverge on the axis. The
        quantum smoothing is a Bessel function $J_0^2(\dots)$.
  \item \textbf{Orbiting.} If $E$ equals a maximum of the effective potential
        $V+L^2/(2\mu r^2)$, $\Theta\to-\infty$ as $b\to b_c$: the particle
        circles the centre many times.
\end{itemize}
Potentials in which these effects occur are typically finite-range and are not
treated here. For the Coulomb case $\Theta(b)$ is monotonic, there is a single
branch, and \eqref{eq:xs-sum} applies directly.

\section{Rutherford scattering}

For $V=\alpha/r$ put $u=1/r$. The orbit equation $u''+u=-\mu\alpha/L^2$ gives
conic sections, here hyperbolas with eccentricity
\[
  \varepsilon=\sqrt{1+(2Eb/\alpha)^2}.
\]
The asymptotes are the zeros of $u$. For repulsion
$u=(\varepsilon\cos\varphi'-1)/p$ with $p=L^2/(\mu\alpha)$, so
$\cos(\Delta\varphi/2)=1/\varepsilon$, $\tan(\Delta\varphi/2)=2Eb/\alpha$, and
with $\Theta=\pi-\Delta\varphi$,
\begin{equation}\label{eq:b-ruth}
  b=\frac{\alpha}{2E}\cot\frac\Theta2 .
\end{equation}
The attractive case gives the same relation with $\alpha$ replaced by
$|\alpha|$ and $\Theta\to-\Theta$. Substituting in \eqref{eq:xs-sum} with
$\dd b/\dd\theta=-(\alpha/4E)/\sin^2(\theta/2)$,
\begin{equation}\label{eq:ruth}
  \dsdO=\left(\frac{\alpha}{4E}\right)^{2}\frac1{\sin^4(\theta/2)} .
\end{equation}
This is Rutherford's formula. Notes:
\begin{itemize}
  \item It depends on $\alpha^2$ and therefore does not distinguish attraction
        from repulsion.
  \item It is exactly proportional to $E^{-2}$ and $\alpha^2$, so one scale
        ($\alpha/E$, a length) fixes the whole angular distribution.
  \item It is singular as $\theta\to0$ like $\theta^{-4}$. The total cross
        section is infinite, because a particle at any impact parameter,
        however large, is deflected a little ($\Theta\approx\alpha/Eb$ for
        large $b$), and the area $\int2\pi b\,\dd b$ is unbounded.
\end{itemize}

\section{Energy transfer and the Coulomb logarithm}\label{sec:coullog}

For a heavy target $m_2$ at rest (see \eqref{eq:recoil}), the recoil energy is
$T=T_{\max}\sin^2(\theta^*/2)$. Changing variables in \eqref{eq:ruth} gives
\begin{equation}\label{eq:dsdT}
  \frac{\dd\sigma}{\dd T}=\frac{2\pi\alpha^2}{m_2v^2\,T^2},
\end{equation}
where $v$ is the projectile's lab speed. Notice that the result depends on the
projectile only through its speed. The energy lost by a fast particle per unit
path length in a medium with $n_2$ targets per unit volume,
$\dd E/\dd x=n_2\!\int T\,\dd\sigma$, then diverges logarithmically at small
$T$, since $\int\dd T/T=\ln(T_{\max}/T_{\min})$. The same logarithm appears in
the momentum-transfer cross section
\begin{equation}\label{eq:sigtr}
  \sigma_{\mathrm{tr}}=\int(1-\cos\theta)\,\dd\sigma
  =\pi\Bigl(\frac\alpha E\Bigr)^{2}\ln\frac{1}{\sin(\theta_{\min}/2)} .
\end{equation}
The cutoff $\theta_{\min}$ (or $T_{\min}$) is supplied by something outside the
pure Coulomb problem, such as screening by other charges. This is the Coulomb
logarithm of plasma physics and stopping-power theory; its appearance is the
practical face of the infinite range.

\section{Infinite range: scaling and the divergent total cross section}
\label{sec:scaling}

For $V=\alpha/r^n$ the only length that can be formed from $\alpha$ and $E$ is
$(\alpha/E)^{1/n}$, so $b=(\alpha/E)^{1/n}g_n(\Theta)$ and
\begin{equation}\label{eq:scaling}
  \dsdO=\Bigl(\frac\alpha E\Bigr)^{2/n}F_n(\theta).
\end{equation}
For large $b$ the deflection decays as $\Theta\sim b^{-n}$, so the cross section
for deflection by more than $\theta_{\min}$ grows like
$\theta_{\min}^{-2/n}$ and diverges as $\theta_{\min}\to0$.
\emph{The classical total cross section is infinite for every infinite-range
potential}, however rapidly it falls off. As we shall see in Chapter~4, quantum
mechanics is more forgiving: diffraction makes very small deflections
unobservable, and the quantum total cross section is finite if and only if $V$
falls faster than $r^{-2}$ (the Coulomb case, $n=1$, remains infinite).

\section{The inverse-square potential}\label{sec:classinv2}

The case $V=\alpha/r^2$ is exactly solvable and shows the generic structure.
The radial equation is that of a free particle with an effective angular
momentum, $L'^2=L^2+2\mu\alpha$. A free particle sweeps an angle $\pi$ in the
variable $\varphi'=\int L'\,\dd t/(\mu r^2)$, so the physical angle is
$\Delta\varphi=\pi L/L'$, and
\begin{equation}\label{eq:Theta2}
  \Theta(b)=\pi\left[1-\frac{b}{\sqrt{b^2+\alpha/E}}\right].
\end{equation}
Inverting,
$b^2=(\alpha/E)(\pi-\Theta)^2/[\Theta(2\pi-\Theta)]$, and from
\eqref{eq:xs-sum}
\begin{equation}\label{eq:xs2}
  \dsdO=\frac{\alpha}{E}\;
   \frac{\pi^2(\pi-\theta)}{\theta^2(2\pi-\theta)^2\sin\theta}.
\end{equation}
At small angles $\dd\sigma/\dd\Omega\approx(\pi\alpha/4E)/\theta^3$, a lower
power of $1/\theta$ than Rutherford's $\theta^{-4}$, but the total cross
section still diverges. For attraction ($\alpha<0$) two new features arise. If
$b^2<|\alpha|/E$, then $L'^2<0$ and the particle spirals into the centre; the
capture cross section is
\begin{equation}\label{eq:cap2}
  \sigma_{\mathrm{cap}}=\frac{\pi|\alpha|}{E}.
\end{equation}
And as $b\to\sqrt{|\alpha|/E}$ from above, $\Theta\to-\infty$: orbiting. In both
cases the unique scale is $\alpha/E$, and $\sigma_{\mathrm{cap}}\propto1/E$. We
shall meet the quantum version of this problem in Section~\ref{sec:qinv2},
where the semiclassical relation turns out to be exact.

\section{Problems}
\begin{enumerate}
  \item Derive \eqref{eq:b-ruth} directly from \eqref{eq:phi-swept} by
        evaluating the integral for $V=\alpha/r$.
  \item Show that for the attractive $1/r^2$ potential the deflection function
        \eqref{eq:Theta2}, continued to $\alpha<0$, gives infinitely many
        branches for $\theta$ and compute the cross section as a sum over
        windings $n$.
  \item For electrons of 100~keV scattering from gold ($Z=79$), find the impact
        parameter that gives $\theta=90^\circ$ and the classical cross section
        per unit solid angle at that angle, and explain why the nucleus can be
        treated as a point.
\end{enumerate}

\section*{Notes and references}
The classical theory of scattering in a central field is in \cite{LL1,Goldstein}; Rutherford's original analysis is \cite{Rutherford11}. The classification of rainbow, glory and orbiting singularities and their quantum smoothing is treated in \cite{FordWheeler,BerryMount}. The Coulomb logarithm and the energy-transfer cross section are discussed in the context of stopping power in \cite{Jackson}. The inverse-square potential and the fall to the centre are treated in \cite{LL3,Case50}; the underlying questions of self-adjointness are explained in \cite{ReedSimonII}.

% =====================================================================
\part{Quantum Scattering}
% =====================================================================

\chapter{Quantum Scattering Without Finite Range}

\section{Stationary scattering states and the cross section}

A scattering experiment is a time-dependent process: a wave packet of mean
momentum $\hbar k$, prepared far from the target, approaches, interacts and
separates. For a packet that is sharp in momentum and broad in space, the
description reduces to a stationary one: a solution of
$(-\hbar^2\nabla^2/2\mu+V)\psi=E\psi$, with $E=\hbar^2k^2/2\mu$, that behaves
at large $r$ as
\begin{equation}\label{eq:asym}
  \psi_{\vk}(\vr)\to \ee^{\ii\vk\cdot\vr}+f(\theta,\varphi)\,\frac{\ee^{\ii kr}}{r}.
\end{equation}
The incident plane wave carries current density $\hbar k/\mu$. The outgoing
spherical wave carries a radial current $(\hbar k/\mu)|f|^2/r^2$, so the flux
through $r^2\dd\Omega$ divided by the incident flux gives
\begin{equation}\label{eq:xs-f}
  \dsdO=|f(\theta,\varphi)|^2 .
\end{equation}

The form \eqref{eq:asym} presupposes that the particle moves freely at large
distances, and the Coulomb potential sits exactly on the borderline where this
fails. In the WKB picture the local wave number is
$k(r)\approx k-\mu V(r)/(\hbar^2k)$, so the phase accumulated out to radius $r$
contains $-(\mu/\hbar^2k)\int^{r}V\,\dd r'$. If $V$ falls faster than $1/r$ this
integral converges and the asymptotic wave is a plane wave. If $V\propto1/r$ it
grows like $\ln r$, so the incident wave carries a logarithmic phase and the
outgoing wave is $\ee^{\ii(kr-\eta\ln2kr)}/r$, not $\ee^{\ii kr}/r$. The
standard form \eqref{eq:asym} therefore holds for potentials that decay faster
than $1/r$; a finite total cross section requires a still faster decay, as we
saw in Section~\ref{sec:scaling} and will quantify below.

\section{Formal structure: the T-matrix, unitarity and reciprocity}

The stationary state with outgoing-wave boundary conditions obeys the
Lippmann--Schwinger equation
\begin{equation}\label{eq:LS}
  \psi^+_{\vk}(\vr)=\ee^{\ii\vk\cdot\vr}+\int G_0^+(\vr,\vr')\,V(\vr')\,
  \psi^+_{\vk}(\vr')\,\dd^3r',\qquad
  G_0^+(\vr,\vr')=-\frac{\mu}{2\pi\hbar^2}\frac{\ee^{\ii k|\vr-\vr'|}}{|\vr-\vr'|}.
\end{equation}
Taking $r\to\infty$ along the direction $\vk'=k\hat{\vr}$ identifies the
amplitude exactly:
\begin{equation}\label{eq:fT}
  f(\vk',\vk)=-\frac{\mu}{2\pi\hbar^2}\,\langle\vk'|V|\psi^+_{\vk}\rangle
  =-\frac{\mu}{2\pi\hbar^2}\,\langle\vk'|T(E+\ii0)|\vk\rangle ,
\end{equation}
with $\langle\vr|\vk\rangle=\ee^{\ii\vk\cdot\vr}$ and $T=V+VG_0T$. We use
\eqref{eq:LS}--\eqref{eq:fT} only as definitions; they are never expanded in
powers of $V$. Several exact properties follow.
\begin{itemize}
  \item \textbf{Unitarity.} Probability conservation makes the S-matrix unitary.
        On the energy shell, for any interaction for which the amplitude exists
        and all channels are elastic,
        \begin{equation}\label{eq:unitarity}
          f(\vk',\vk)-f^*(\vk,\vk')=\frac{\ii k}{2\pi}\int\dd\Omega''\,
          f^*(\vk'',\vk')\,f(\vk'',\vk).
        \end{equation}
  \item \textbf{Optical theorem.} Setting $\vk'=\vk$ in \eqref{eq:unitarity},
        \begin{equation}\label{eq:optical}
          \Imag f(\vk,\vk)=\frac{k}{4\pi}\,\sigma_{\mathrm{tot}} .
        \end{equation}
        The forward amplitude is the interference of the scattered wave with
        the incident wave, and its imaginary part is the shadow cast by
        whatever is removed from the beam. If inelastic channels are open,
        \eqref{eq:optical} still holds with
        $\sigma_{\mathrm{tot}}=\sigma_{\mathrm{el}}+\sigma_{\mathrm{inel}}$,
        while \eqref{eq:unitarity} becomes an inequality.
  \item \textbf{Reciprocity.} Time-reversal invariance gives
        $f(\vk',\vk)=f(-\vk,-\vk')$. For a central, real potential $f$ depends
        only on the angle $\theta$ between $\vk$ and $\vk'$, and
        $f(\vk',\vk)=f(\vk,\vk')$.
  \item \textbf{Coulomb.} For $V=\alpha/r$ the integral in \eqref{eq:fT} is only
        conditionally defined and the forward amplitude is infinite. This is
        consistent with \eqref{eq:optical}, since $\sigma_{\mathrm{tot}}$ is
        infinite, but it means that \eqref{eq:optical} cannot be used as a
        tool for the pure Coulomb problem.
\end{itemize}

\section{Long-range forces and Dollard's asymptotic states}\label{sec:dollard}

The physical content of a scattering state is that the true evolution
$\ee^{-\ii Ht/\hbar}\psi$ resembles a freely moving packet as $t\to\mp\infty$.
Standard scattering theory compares $H=H_0+V$ with $H_0$ through the
M{\o}ller operators
$\Omega_\pm=\lim\,\ee^{\ii Ht/\hbar}\ee^{-\ii H_0t/\hbar}$ as $t\to\mp\infty$.
For $V=\alpha/r$ this limit does not exist: a packet moving at speed $v$ is at
distance $r=vt$, and the phase accumulated from the potential grows without
bound as $-\eta(p)\ln|t|$. Dollard (1964) repaired the theory by modifying the
comparison dynamics:
\begin{equation}\label{eq:dollard}
  U_D(t)=\exp\Bigl\{-\frac{\ii}{\hbar}H_0t-\ii\,\eta(\hat p)\,\sgn(t)\,
  \ln\frac{4H_0|t|}{\hbar}\Bigr\},\qquad \eta(p)=\frac{\mu\alpha}{\hbar p}.
\end{equation}
The modified wave operators
$\Omega^D_\pm=\lim\,\ee^{\ii Ht/\hbar}U_D(t)$ as $t\to\mp\infty$ exist, are
isometric and complete, and define a unitary
$S=(\Omega^D_+)^\dagger\Omega^D_-$, whose kernel is the Coulomb amplitude. One
consequence is conceptually important. The scale inside the logarithm is
arbitrary: replacing $4H_0|t|/\hbar$ by $|t|/\tau$ multiplies $S$ by a phase
that depends on the energy only. Hence:
\begin{itemize}
  \item $|f_C(\theta)|^2$, the Rutherford cross section, is unambiguous;
  \item the overall phase of $f_C$ is a matter of convention;
  \item the phase of one Coulomb amplitude relative to another at the same
        energy (as when two amplitudes interfere) is physical.
\end{itemize}
A real beam is screened at some distance $R$ by other charges, so that ordinary
theory applies. As $R\to\infty$ the screened amplitude approaches $f_C$ times a
phase factor that depends only on $k$ and $R$ and diverges logarithmically; the
cross section at any fixed $\theta\neq0$ converges to Rutherford's.

\section{Problems}
\begin{enumerate}
  \item Show by a flux argument that \eqref{eq:unitarity} with $\vk'=\vk$ leads
        to \eqref{eq:optical}.
  \item Verify that for $V\propto r^{-n}$ with $n>1$ the accumulated WKB phase
        converges, and find the leading large-$r$ correction to the plane wave
        for $1<n<2$.
  \item Show that changing the scale of the logarithm in \eqref{eq:dollard}
        shifts the phase of $f_C$ by a $\theta$-independent amount and leaves
        $|f_C|^2$ unchanged.
\end{enumerate}

\section*{Notes and references}
General formal scattering theory, including the Lippmann--Schwinger equation \cite{LS50}, the T-matrix, unitarity and the optical theorem, can be found in \cite{Newton,Taylor,GW,Messiah}. The failure of the standard asymptotic condition for the Coulomb potential and its repair were given by Dollard \cite{Dollard64}; a complete mathematical treatment of long-range scattering theory is in \cite{ReedSimonIII}.

% =====================================================================
\chapter{The Exact Coulomb Problem}
% =====================================================================

\section{Separation in parabolic coordinates}

For $V=\alpha/r$ the Schr\"odinger equation is
\begin{equation}\label{eq:coul-schr}
  \Bigl[\nabla^2+k^2-\frac{2k\eta}{r}\Bigr]\psi=0,\qquad
  \eta=\frac{\mu\alpha}{\hbar^2k}=\frac{\alpha}{\hbar v}.
\end{equation}
The equation separates in parabolic coordinates $\xi=r-z$, $\chi=r+z$,
$\varphi$, where the $z$-axis is the incident direction. The Laplacian is
\[
  \nabla^2=\frac{4}{\xi+\chi}\Bigl[\partial_\xi(\xi\partial_\xi)
  +\partial_\chi(\chi\partial_\chi)\Bigr]+\frac1{\xi\chi}\partial_\varphi^2 .
\]
For an axially symmetric solution of the form $\psi=\ee^{\ii kz}w(\xi)$, which
carries a plane-wave factor along the beam, the $\chi$-dependence cancels and
\begin{equation}\label{eq:kummer-w}
  \xi\,w''+(1-\ii k\xi)\,w'-k\eta\,w=0 .
\end{equation}
With $s=\ii k\xi$ this is Kummer's equation
$s\,w_{ss}+(1-s)\,w_s-a\,w=0$ with $a=-\ii\eta$ and $b=1$. Its second solution
is logarithmically singular on the forward axis $\xi=0$, which is unphysical,
so
\begin{equation}\label{eq:psiC}
  \psi_C(\vr)=C\,\ee^{\ii kz}\,\Fhyp\bigl(-\ii\eta;1;\ii k(r-z)\bigr).
\end{equation}

\section{Asymptotics and the scattering amplitude}

For large $|s|$,
$\Fhyp(a;1;s)\simeq(-s)^{-a}/\Gamma(1-a)+\ee^{s}s^{a-1}/\Gamma(a)$. Put
$a=-\ii\eta$, $s=\ii k\xi$, $\xi=r(1-\cos\theta)=2r\sin^2(\theta/2)$. The two
terms are
\begin{align*}
  \frac{(-\ii k\xi)^{\ii\eta}}{\Gamma(1+\ii\eta)}
    &=\frac{\ee^{\pi\eta/2}\,\ee^{\ii\eta\ln k\xi}}{\Gamma(1+\ii\eta)},\\[2pt]
  \frac{\ee^{\ii k\xi}(\ii k\xi)^{-\ii\eta-1}}{\Gamma(-\ii\eta)}
    &=\frac{\ee^{\pi\eta/2}\,\ee^{\ii k\xi-\ii\eta\ln k\xi}}
           {\ii k\xi\,\Gamma(-\ii\eta)} .
\end{align*}
Choosing $C=\ee^{-\pi\eta/2}\Gamma(1+\ii\eta)$ normalises the first term to unit
amplitude, and using
$\Gamma(1+\ii\eta)/\Gamma(-\ii\eta)=-\ii\eta\,\Gamma(1+\ii\eta)/\Gamma(1-\ii\eta)$
for the second,
\begin{align}
  \psi_C&\to \ee^{\ii[kz+\eta\ln k(r-z)]}
        +f_C(\theta)\,\frac{\ee^{\ii[kr-\eta\ln2kr]}}{r}, \label{eq:psiC-asym}\\[4pt]
  f_C(\theta)&=-\frac{\eta}{2k\sin^2(\theta/2)}\,
        \exp\!\bigl[-\ii\eta\ln\sin^2(\theta/2)+2\ii\sigma_0\bigr],
        \qquad \sigma_0=\arg\Gamma(1+\ii\eta). \label{eq:fC}
\end{align}
The result holds for $kr(1-\cos\theta)\gg1$, that is away from the forward axis,
and for either sign of $\eta$.

\section{Rutherford's law, exactly}

Since $\eta/k=\alpha/2E$,
\begin{equation}\label{eq:ruth-q}
  \dsdO=|f_C|^2=\Bigl(\frac{\alpha}{4E}\Bigr)^{2}\frac1{\sin^4(\theta/2)} .
\end{equation}
This is the classical result \eqref{eq:ruth}, and $\hbar$ has dropped out
completely. The agreement is exact and holds for every value of $\eta$, large
or small.

\section{Anatomy of the Coulomb amplitude}

\paragraph{The incident wave is distorted.} The unscattered wave in
\eqref{eq:psiC-asym} is a plane wave with the extra phase $\eta\ln k(r-z)$.
There is no region of space in which the particle is simply free, and the
forward direction cannot be cleanly separated from the beam. The amplitude
\eqref{eq:fC} diverges as $\theta^{-2}$ as $\theta\to0$, and
\eqref{eq:psiC-asym} fails inside the paraboloid $r-z\lesssim1/k$ around the
axis.

\paragraph{The phase of $f_C$.} The phase $2\sigma_0-\eta\ln\sin^2(\theta/2)$
depends on $\theta$ and is observable only through interference, for example
between the two amplitudes for identical particles (Section~\ref{sec:mott}).

\paragraph{Bound states are poles.} Since
$\ee^{2\ii\sigma_0}=\Gamma(1+\ii\eta)/\Gamma(1-\ii\eta)$, the amplitude has
poles where $1+\ii\eta=-m$, $m=0,1,2,\dots$, that is at $\eta=\ii n$ with
$n=m+1$. For an attractive potential, $\alpha=-|\alpha|$, continue to imaginary
momentum $k=\ii\kappa$: $\eta=-\mu|\alpha|/(\hbar^2\ii\kappa)
=\ii\mu|\alpha|/(\hbar^2\kappa)$. The pole condition gives
$\kappa=\mu|\alpha|/(\hbar^2n)$, hence
\begin{equation}\label{eq:bohr}
  E_n=-\frac{\mu\alpha^2}{2\hbar^2n^2},\qquad n=1,2,3,\dots
\end{equation}
The complete Bohr spectrum is encoded in the scattering amplitude, a prototype
of the general statement that bound states are poles of the S-matrix at
negative energy.

\paragraph{Penetration factor.} The probability density at the origin is
$|\psi_C(0)|^2=|C|^2=\ee^{-\pi\eta}|\Gamma(1+\ii\eta)|^2$, and using
$|\Gamma(1+\ii\eta)|^2=\pi\eta/\sinh\pi\eta$,
\begin{equation}\label{eq:sommerfeld}
  |\psi_C(0)|^2=\frac{2\pi\eta}{\ee^{2\pi\eta}-1}.
\end{equation}
For repulsion with $\eta\gg1$ this is $\ee^{-2\pi\eta}$, the Gamow tunnelling
factor that governs nuclear fusion at low energy. For attraction with
$|\eta|\gg1$ it is $\approx2\pi|\eta|$, the Sommerfeld enhancement of reactions
that happen at the origin, relevant for threshold production and annihilation of
charged pairs. As $\eta\to0$ the factor tends to~1.

\section{Partial waves and the semiclassical dictionary}\label{sec:partial}

Even though the potential has infinite range, the problem can be expanded in
partial waves, in terms of Coulomb functions $F_\ell(\eta,\rho)$ that solve
$u''+[k^2-2k\eta/r-\ell(\ell+1)/r^2]u=0$:
\begin{equation}\label{eq:cwaves}
  \begin{gathered}
  \psi_C=\sum_\ell(2\ell+1)\,\ii^\ell\,\ee^{\ii\sigma_\ell}\,
  \frac{F_\ell(\eta,kr)}{kr}\,P_\ell(\cos\theta),\\[4pt]
  F_\ell\to\sin\!\bigl(\rho-\eta\ln2\rho-\tfrac{\ell\pi}{2}+\sigma_\ell\bigr),
  \qquad \sigma_\ell=\arg\Gamma(\ell+1+\ii\eta).
  \end{gathered}
\end{equation}
The amplitude is then
\begin{equation}\label{eq:fC-pw}
  f_C(\theta)=\frac1{2\ii k}\sum_\ell(2\ell+1)
  \bigl(\ee^{2\ii\sigma_\ell}-1\bigr)P_\ell(\cos\theta),
\end{equation}
where the $-1$ contributes only at $\theta=0$, because
$\sum(2\ell+1)P_\ell=2\delta(1-\cos\theta)$. The phases obey
$\sigma_\ell=\sigma_{\ell-1}+\arctan(\eta/\ell)$. Two warnings distinguish this
from the finite-range case. First, the terms in \eqref{eq:fC-pw} do not decay:
their modulus is $(2\ell+1)/2k$ at large $\ell$. The series converges only as a
distribution, for example by Abel summation, and its sum is \eqref{eq:fC}.
Second, the phases $\sigma_\ell$ do not tend to a constant but grow like
$\eta\ln\ell$, which is the partial-wave form of the logarithmic phase of
Section~\ref{sec:dollard}.

The semiclassical dictionary connects the phases to classical deflection. In
the convention where positive $\Theta$ means repulsion, the WKB phase shift
obeys $\Theta(\ell)=2\,\dd\delta_\ell/\dd\ell$, with $\ell+\tfrac12=kb$. Applied
to the Coulomb phases,
\[
  \frac{\dd\sigma_\ell}{\dd\ell}=\Imag\,\psi(\ell+1+\ii\eta)
  \simeq\Imag\ln\bigl(\ell+\tfrac12+\ii\eta\bigr)
  =\arctan\frac{\eta}{\ell+\tfrac12},
\]
where $\psi$ is the digamma function and the approximation is accurate to
$O(1/\ell^2)$. Therefore
\begin{equation}\label{eq:dict}
  \tan\frac\Theta2=\frac{\eta}{\ell+\tfrac12}=\frac{\alpha}{2Eb}.
\end{equation}
This is exactly the classical deflection function \eqref{eq:b-ruth}. The
Coulomb problem is the case where the dictionary is exact.

\section{Identical particles: the Mott formula}\label{sec:mott}

For two identical particles the exchange $\vr\to-\vr$ means
$\theta\to\pi-\theta$, and the amplitude must be symmetrised. Write
$s=\sin(\theta/2)$, $c=\cos(\theta/2)$, and use the Coulomb amplitude
\eqref{eq:fC} with the CM energy $E$ and reduced mass $\mu=m/2$. Remember that
$\eta/2k=\alpha/4E$.
\begin{itemize}
  \item \textbf{Spin 0} (or spin-polarised bosons): the cross section is
        $|f(\theta)+f(\pi-\theta)|^2$,
        \begin{equation}\label{eq:mott0}
          \dsdO=\Bigl(\frac{\alpha}{4E}\Bigr)^{2}\Bigl[\frac1{s^4}+\frac1{c^4}
          +\frac{2}{s^2c^2}\cos\bigl(\eta\ln\tan^2\tfrac\theta2\bigr)\Bigr].
        \end{equation}
  \item \textbf{Spin $\tfrac12$, unpolarised:} the singlet (weight $\tfrac14$) is
        spatially symmetric and the triplet (weight $\tfrac34$) antisymmetric, so
        $\dd\sigma/\dd\Omega=|f|^2+|f'|^2-\Real(ff'^*)$ with
        $f'=f(\pi-\theta)$:
        \begin{equation}\label{eq:mott12}
          \dsdO=\Bigl(\frac{\alpha}{4E}\Bigr)^{2}\Bigl[\frac1{s^4}+\frac1{c^4}
          -\frac{1}{s^2c^2}\cos\bigl(\eta\ln\tan^2\tfrac\theta2\bigr)\Bigr].
        \end{equation}
\end{itemize}
At $\theta=90^\circ$ these give $16$, $4$ and, for distinguishable particles,
$8$ in units of $(\alpha/4E)^2$: bosonic doubling and fermionic suppression of
the classical value. The cosine terms carry the logarithmic Coulomb phase
$\eta\ln\tan^2(\theta/2)$, so the unobservable-looking phase of
Section~\ref{sec:dollard} shows up in the angular distribution. For $\eta\gg1$
the cosine oscillates rapidly with $\theta$ and averages to zero over any
realistic detector resolution, leaving the classical sum of Rutherford terms:
quantum interference gives way to classical addition of probabilities.

\section{The inverse-square potential in quantum mechanics}\label{sec:qinv2}

For $V=\alpha/r^2$ the radial equation is
$u''+[k^2-(\ell(\ell+1)+g)/r^2]u=0$ with $g=2\mu\alpha/\hbar^2$. This is the
free radial equation with $\ell$ replaced by a non-integer $\lambda$, where
$\lambda+\tfrac12=\sqrt{(\ell+\tfrac12)^2+g}$. The regular solutions are Bessel
functions of order $\lambda+\tfrac12$, so
\begin{equation}\label{eq:delta2}
  \delta_\ell=\frac\pi2\Bigl[\ell+\tfrac12-\sqrt{(\ell+\tfrac12)^2+g}\Bigr].
\end{equation}
The phase shift is independent of the energy, a reflection of scale invariance.
Three facts follow.
\begin{itemize}
  \item Applying $\Theta=2\,\dd\delta_\ell/\dd\ell$ with $kb=\ell+\tfrac12$ and
        $g/k^2=\alpha/E$, one finds
        $\Theta=\pi\bigl[1-b/\sqrt{b^2+\alpha/E}\bigr]$, exactly the classical
        deflection \eqref{eq:Theta2}.
  \item At large $\ell$, $\delta_\ell\approx-\pi g/4\ell$, so
        $\sum(2\ell+1)\sin^2\delta_\ell$ diverges logarithmically: the total
        cross section is infinite. This is the boundary case $n=2$ of the rule
        stated in Section~\ref{sec:scaling}, that the total cross section is
        finite if and only if $V$ falls faster than $r^{-2}$.
  \item For attraction, if $g<-(\ell+\tfrac12)^2$ then $\lambda+\tfrac12$ is
        imaginary. This is the quantum fall to the centre, and the condition is
        the same as the classical capture condition $b^2<|\alpha|/E$, since
        $(\ell+\tfrac12)^2=k^2b^2$. The Hamiltonian then fails to be
        self-adjoint, and the physics requires a self-adjoint extension with a
        new length scale.
\end{itemize}

\section{When is the classical description valid?}

A wave packet localised to an impact parameter $b$ has an angular uncertainty
$\delta\theta\approx1/(kb)$, while the classical deflection is
$\Theta\approx2\eta/(kb)$. Classical trajectories are meaningful when
$\Theta\gg\delta\theta$, that is when $\eta\gg1$. For the Coulomb potential
$\eta=Z_1Z_2\,(e^2/\hbar c)\,(c/v)$, so heavy ions at low speed are classical
($\eta$ of order 10--100), while fast electrons on light nuclei are in the wave
regime. For the inverse-square potential the criterion is $g\gg1$.

It is striking that the exact Coulomb cross section \eqref{eq:ruth-q} coincides
with the classical one even for $\eta\ll1$, where the classical picture has no
meaning. This is a property of the $1/r$ potential alone, tied to its hidden
dynamical symmetry ($SO(4)$ for bound states, $SO(3,1)$ for scattering,
generated by the Runge--Lenz vector). It fails for the screened Coulomb
potential, it fails for $1/r^2$ (where the deflection function is exact but the
quantum cross section differs from \eqref{eq:xs2} by interference), and, as
Part~III shows, it fails for relativistic motion.

\section{Problems}
\begin{enumerate}
  \item Carry out the substitution that leads from \eqref{eq:coul-schr} to
        \eqref{eq:kummer-w} using the parabolic form of $\nabla^2$.
  \item Prove $\sigma_\ell-\sigma_{\ell-1}=\arctan(\eta/\ell)$ from
        $\Gamma(z+1)=z\Gamma(z)$, and use it to show $\sigma_\ell\simeq\eta\ln\ell$
        at large $\ell$.
  \item Derive \eqref{eq:mott12} from the singlet and triplet weights, and
        evaluate the cross section for $\alpha$--$\alpha$ scattering at
        $90^\circ$ for $\eta=0.5$ and $\eta=5$, treating the $\alpha$ particles
        as spin~0.
  \item Use \eqref{eq:sommerfeld} to estimate the Gamow factor for $p+p$ at
        1~keV CM energy, and explain why fusion in the Sun, at about 1~keV,
        proceeds only through the high-energy tail of the thermal distribution.
\end{enumerate}

\section*{Notes and references}
The exact Coulomb scattering amplitude was obtained independently by Gordon \cite{GordonB} and Mott \cite{Mott28}. Textbook treatments are in \cite{LL3,Messiah,Newton}. Coulomb and Kummer functions are tabulated in \cite{AS,DLMF}, with a modern numerical treatment in \cite{Seaton02}; useful integrals are in \cite{GR}. The penetration factor of Section~4.4 goes back to Gamow \cite{Gamow28} and Sommerfeld \cite{Sommerfeld31}. The hidden $SO(4)$ symmetry of the Coulomb problem was found by Pauli, Fock and Bargmann \cite{Pauli26,Fock35,Bargmann36} and is reviewed in \cite{Bander66}. The scattering of identical particles by the Coulomb force was treated by Mott \cite{Mott30}.

% =====================================================================
\part{Relativistic Scattering}
% =====================================================================

\chapter{Classical Relativistic Scattering}

The relativistic two-body problem has no simple Hamiltonian formulation with
instantaneous interactions, so we treat a test particle of mass $m$ in the
static field of an infinitely heavy source. This is a legitimate limit for
$m\ll M$ and the classical starting point for the Klein--Gordon and Dirac
problems of Chapter~6. Radiation losses are assumed negligible.

\section{The relativistic Coulomb orbit}

Let $E$ be the total energy, $p$ the momentum at infinity (so
$E^2=p^2c^2+m^2c^4$), $v=pc^2/E$ the speed at infinity, and $L=pb$. Energy
conservation reads $(E-\alpha/r)^2=p_r^2c^2+L^2c^2/r^2+m^2c^4$. With $u=1/r$
this gives
\begin{equation}\label{eq:pr2}
  p_r^2=p^2-\frac{2E\alpha}{c^2}\,u-\Bigl(L^2-\frac{\alpha^2}{c^2}\Bigr)u^2 .
\end{equation}
The Hamilton--Jacobi equation gives $\dd\varphi=L\,\dd r/(r^2p_r)$. Define
\begin{equation}\label{eq:kappa}
  \kappa=\sqrt{1-\frac{\alpha^2}{L^2c^2}}=\sqrt{1-\Bigl(\frac{\alpha}{pcb}\Bigr)^2}.
\end{equation}
The integral is elementary and the orbit is
$u(\varphi)=u_c+D\cos[\kappa(\varphi-\varphi_0)]$ with
$u_c=-E\alpha/(c^2L^2\kappa^2)$. Compare the Newtonian conic, which has
$\kappa=1$: the relativistic orbit is a precessing rosette whose angular period
is $2\pi/\kappa$ rather than $2\pi$.

\section{The deflection function}

The asymptotes are the zeros of $u$. Solving $u(\varphi)=0$ gives, for either
sign of $\alpha$ and $L>|\alpha|/c$, the swept angle and deflection
\begin{equation}\label{eq:dphi-rel}
  \Delta\varphi=\frac2\kappa\Bigl[\frac\pi2-\arctan\frac{\alpha}{\kappa\,b\,p\,v}\Bigr],
  \qquad\Theta=\pi-\Delta\varphi .
\end{equation}
For $\kappa\to1$ and $pv\to2E$ this is $\tan(\Theta/2)=\alpha/(2Eb)$,
Rutherford's result \eqref{eq:b-ruth}. For weak coupling,
\begin{equation}\label{eq:theta-exp}
  \Theta=\frac{2\alpha}{pvb}-\frac{\pi\alpha^2}{2p^2c^2b^2}+O(\alpha^3).
\end{equation}
The first term is the familiar impulse approximation, with $2E$ replaced by
$pv$. The second term is a relativistic effect with no non-relativistic
analogue: it comes from the precession factor $\kappa$ and is the scattering
counterpart of the perihelion advance. Its size relative to the first term is
$(\pi/4)(v/c)(\alpha/pcb)$.

Because \eqref{eq:dphi-rel} is transcendental in $b$, the cross section
\eqref{eq:xs-sum} has no closed form. To leading order in $\alpha$,
$\tan(\Theta/2)=\alpha/(pvb)$ and
\begin{equation}\label{eq:ruth-rel}
  \dsdO=\Bigl(\frac{\alpha}{2pv}\Bigr)^{2}\frac1{\sin^4(\theta/2)},
\end{equation}
the relativistic Rutherford formula, in which the kinetic energy $2E$ of the
non-relativistic law is replaced by $pv$.

\section{Capture and the fall to the centre}\label{sec:capture}

The coefficient of $u^2$ in \eqref{eq:pr2} changes sign when $L<|\alpha|/c$. For
attraction the particle then spirals into the centre, since $p_r^2$ grows
without bound as $r\to0$. The capture cross section is
\begin{equation}\label{eq:cap-rel}
  \sigma_{\mathrm{cap}}=\frac{\pi\alpha^2}{p^2c^2},
\end{equation}
with no counterpart in non-relativistic Coulomb motion. For repulsion with
$L<\alpha/c$, $\kappa$ is imaginary and \eqref{eq:dphi-rel} continues
analytically ($\arctan$ becomes $\artanh$); the particle is always turned back.
For a head-on collision ($L=0$) the turning point is where the kinetic energy
vanishes, $E-\alpha/r=mc^2$, so $r_{\min}=\alpha/(E-mc^2)$, the same as the
non-relativistic $\alpha/T$.

The capture threshold $L_c=\alpha/c$ has a quantum-mechanical meaning:
$L_c/\hbar=\alpha/\hbar c=Z\afs$, with $\afs$ the fine-structure constant. For
$Z<137$ this is below $\hbar$, smaller than the smallest non-zero orbital
angular momentum. The classical capture region is therefore smaller than a
wavelength and cannot occur as a classical process; Section~\ref{sec:dirac}
shows that the Dirac equation indeed has no fall to the centre for $Z\afs<1$.

\section{Invariant form}

With $t=-4p^2\sin^2(\theta/2)$ and \eqref{eq:dsdt}, the leading-order law
\eqref{eq:ruth-rel} reads
\begin{equation}\label{eq:dsdt-rel}
  \frac{\dd\sigma}{\dd t}=\frac{4\pi\alpha^2}{\beta^2\,t^2},
\end{equation}
where $\beta$ is the speed of the projectile in the rest frame of the source. It
holds exactly in non-relativistic classical and quantum mechanics (with $\beta$
the relative speed) and at leading order in $\alpha$ relativistically for
spinless particles. Spin modifies it (Chapter~6). The $1/t^2$ is the
relativistic seed of Chapter~7.

\section{Scattering of radiation by a free charge: Thomson scattering}

A free charge $e$ of mass $m$ in a plane wave of amplitude $E_0$ and frequency
$\omega$ is accelerated to $a=eE_0\cos(\omega t)/m$, provided
$a_0=eE_0/(mc\omega)\ll1$ so that the motion is non-relativistic. The Larmor
formula gives the radiated power per unit solid angle
$\dd P/\dd\Omega=(e^2a^2/4\pi c^3)\sin^2\Theta_a$, where $\Theta_a$ is the angle
between $\bm a$ and the direction of observation. The time-averaged incident
flux is $S=cE_0^2/8\pi$, so
\begin{equation}\label{eq:thomson-b}
  \dsdO=r_e^2\sin^2\Theta_a,\qquad r_e=\frac{e^2}{mc^2}.
\end{equation}
For an unpolarised beam, averaging over polarisations gives
\begin{equation}\label{eq:thomson-u}
  \dsdO=\frac{r_e^2}{2}\,(1+\cos^2\theta),\qquad
  \sigma_T=\frac{8\pi}{3}\,r_e^2=6.65\times10^{-25}\ \mathrm{cm}^2 .
\end{equation}
The cross section is independent of frequency. It holds while
$\hbar\omega\ll mc^2$, so that photon recoil is negligible; the Klein--Nishina
corrections for larger $\hbar\omega$ come from perturbative quantum
electrodynamics and are omitted here. When $a_0\gtrsim1$ the electron becomes
relativistic in the wave, and the scattering becomes nonlinear: the scattered
light contains harmonics of $\omega$, and the cycle-averaged mass is increased to
$m\sqrt{1+a_0^2/2}$ for linear polarisation ($m\sqrt{1+a_0^2}$ for circular).
This remains exactly solvable for a plane wave.

\section{Problems}
\begin{enumerate}
  \item Derive \eqref{eq:pr2}--\eqref{eq:dphi-rel} from the Hamilton--Jacobi
        equation.
  \item Show that \eqref{eq:dphi-rel} continues to $L<\alpha/c$ for repulsion,
        find the resulting $\Theta(b)$, and verify that $\Theta\to\pi$ as
        $b\to0$.
  \item Show that the relativistic head-on closest approach is
        $r_{\min}=\alpha/(E-mc^2)$ and evaluate it for a 1~MeV electron on
        gold.
  \item Carry out the derivation of \eqref{eq:thomson-b} and
        \eqref{eq:thomson-u} from Larmor's formula.
\end{enumerate}

\section*{Notes and references}
The precession of the relativistic Coulomb orbit is due to Sommerfeld \cite{Sommerfeld16}. The fields of relativistic charges, Thomson scattering and the motion of a charge in a plane wave are treated in \cite{Jackson,LL2}. Thomson scattering is the low-energy limit of the Klein--Nishina result \cite{KN29}, which is perturbative and not treated here.

% =====================================================================
\chapter[Relativistic Quantum Scattering]{Relativistic Quantum Scattering in Exactly Solvable Fields}
% =====================================================================

In this chapter $\zeta=\alpha/\hbar c$ (equal to $Z\afs$ for a nucleus of charge
$Z$) measures how relativistic the Coulomb problem is, while
$\eta=\alpha/\hbar v$ with $v=pc^2/E$ measures how classical it is. The two are
related by $\zeta=\eta v/c$.

\section{The Klein--Gordon equation in a Coulomb field}

For a spin-0 particle in the potential energy $V=\alpha/r$ the Klein--Gordon
equation is $(E-V)^2\psi=(-\hbar^2c^2\nabla^2+m^2c^4)\psi$. Expanding
$(E-\alpha/r)^2-m^2c^4=p^2c^2-2E\alpha/r+\alpha^2/r^2$ gives
\begin{equation}\label{eq:kg}
  \nabla^2\psi+\Bigl[k^2-\frac{2k\eta}{r}+\frac{\zeta^2}{r^2}\Bigr]\psi=0,
  \qquad \hbar k=p,\quad \eta=\frac{E\alpha}{\hbar c^2p}=\frac{\alpha}{\hbar v}.
\end{equation}
For the partial wave $\ell$ the radial equation is
\begin{equation}\label{eq:kg-rad}
  u''+\Bigl[k^2-\frac{2k\eta}{r}-\frac{\ell(\ell+1)-\zeta^2}{r^2}\Bigr]u=0 .
\end{equation}
This is exactly the non-relativistic Coulomb radial equation of
Section~\ref{sec:partial} with $\ell$ replaced by a non-integer $\Lambda$
defined by
\begin{equation}\label{eq:Lambda}
  \Lambda(\Lambda+1)=\ell(\ell+1)-\zeta^2,\qquad
  \Lambda+\tfrac12=\sqrt{(\ell+\tfrac12)^2-\zeta^2}.
\end{equation}
Every result of Chapter~4 therefore carries over, with $\Lambda$ for $\ell$ and
the relativistic $\eta$. The regular solution behaves as
$\sin(kr-\eta\ln2kr-\Lambda\pi/2+\sigma_\Lambda)$, where
$\sigma_\Lambda=\arg\Gamma(\Lambda+1+\ii\eta)$. Measured against the reference
$\sin(kr-\eta\ln2kr-\ell\pi/2)$, the partial-wave phase is
\begin{equation}\label{eq:delta-kg}
  \ee^{2\ii\delta_\ell}=\frac{\Gamma(\Lambda+1+\ii\eta)}{\Gamma(\Lambda+1-\ii\eta)}\,
  \ee^{\ii\pi(\ell-\Lambda)},
\end{equation}
and the amplitude is
$f(\theta)=(2\ii k)^{-1}\sum(2\ell+1)(\ee^{2\ii\delta_\ell}-1)P_\ell(\cos\theta)$,
understood distributionally as in \eqref{eq:fC-pw}. Unlike the non-relativistic
case the series has no known closed-form sum, because the factor
$\ee^{\ii\pi(\ell-\Lambda)}$ spoils the summation that gives \eqref{eq:fC}; the
series is evaluated numerically with convergence acceleration.

\paragraph{Check against the classical orbit.} With $L=\hbar(\ell+\tfrac12)$ one
has $\Lambda+\tfrac12=\kappa(\ell+\tfrac12)$, where $\kappa$ is the precession
factor \eqref{eq:kappa}. The semiclassical rule $\Theta=2\,\dd\delta_\ell/\dd\ell$,
with $\Imag\psi(\Lambda+1+\ii\eta)\simeq\arctan[\eta/(\Lambda+\tfrac12)]$, gives
\[
  \Theta=\pi\Bigl(1-\frac1\kappa\Bigr)+\frac2\kappa\arctan\frac{\alpha}{\kappa\,b\,pv},
\]
which is exactly the classical relativistic deflection function
\eqref{eq:dphi-rel}. At large $\ell$, $\ell-\Lambda\approx\zeta^2/(2\ell+1)$, and
the extra phase in \eqref{eq:delta-kg} reproduces the precession term
$-\pi\alpha^2/(2p^2c^2b^2)$ of \eqref{eq:theta-exp}. The classical relativistic
orbit is thus read off directly from the phases, just as in the non-relativistic
case.

\paragraph{Bound states.} For attraction, $\alpha=-|\alpha|$, the poles of
$\Gamma(\Lambda+1+\ii\eta)$ at $\Lambda+1+\ii\eta=-n_r$ ($n_r=0,1,2,\dots$) give
the energy levels
\begin{equation}\label{eq:kg-levels}
  E=mc^2\Bigl[1+\frac{\zeta^2}{(n_r+\Lambda+1)^2}\Bigr]^{-1/2}.
\end{equation}
These depend on $\ell$ through $\Lambda$, unlike the Dirac levels below.
\emph{Fall to the centre} appears when $\Lambda+\tfrac12$ becomes imaginary, that
is for $\ell+\tfrac12<\zeta$. This is the same condition $L<\alpha/c$ as the
classical capture threshold of Section~\ref{sec:capture}, and for $\ell=0$ it
first occurs at $\zeta=\tfrac12$.

\section{The Dirac equation in a Coulomb field}\label{sec:dirac}

For a spin-$\tfrac12$ particle,
$[c\,\bm\alpha\cdot\bm p+\beta mc^2+\alpha/r]\psi=E\psi$. (Here $\bm\alpha$ and
$\beta$ are the Dirac matrices, not to be confused with the coupling $\alpha$.)
In a central field, $\bm J=\bm L+\bm S$ and parity are conserved, and the radial
problem is labelled by $\kappa_D=-(\ell+1)$ for $j=\ell+\tfrac12$ and
$\kappa_D=+\ell$ for $j=\ell-\tfrac12$, so that $|\kappa_D|=j+\tfrac12$.
Eliminating the small radial component leads again to Coulomb structure, with the
role of $\Lambda+\tfrac12$ played by
\begin{equation}\label{eq:gamma}
  \gamma=\sqrt{\kappa_D^2-\zeta^2}.
\end{equation}
Because $|\kappa_D|\ge1$, $\gamma$ is real for every $\zeta<1$:
\emph{the Dirac Coulomb problem has no fall to the centre for $Z<137$}, in
agreement with the argument of Section~\ref{sec:capture}. The regular solutions
behave as $r^{\gamma-1}$ at the origin, and the bound-state energies are the
Sommerfeld--Dirac levels
\begin{equation}\label{eq:dirac-levels}
  E=mc^2\Bigl[1+\frac{\zeta^2}{(n_r+\gamma)^2}\Bigr]^{-1/2},\qquad
  n_r=n-|\kappa_D| .
\end{equation}
These depend only on $n$ and $j$, so levels with the same $j$ and different
$\ell$ remain degenerate. The scattering phase shifts $\delta_\kappa$ are
likewise given in closed form by ratios of Gamma functions with arguments
$\gamma\pm\ii\eta$ (Mott, 1929), and at large $|\kappa_D|$ one has
$\gamma\approx|\kappa_D|-\zeta^2/(2|\kappa_D|)$, the same structure as for
$\Lambda+\tfrac12$.

A spin-$\tfrac12$ particle scattering from a central field has two independent
amplitudes, one that preserves the spin projection and one that flips it:
\begin{align}
  f(\theta)&=\frac1{2\ii k}\sum_\ell\Bigl[(\ell+1)\bigl(\ee^{2\ii\delta_{-\ell-1}}-1\bigr)
      +\ell\bigl(\ee^{2\ii\delta_\ell}-1\bigr)\Bigr]P_\ell(\cos\theta),\notag\\
  g(\theta)&=\frac1{2\ii k}\sum_\ell\Bigl[\ee^{2\ii\delta_\ell}
      -\ee^{2\ii\delta_{-\ell-1}}\Bigr]P_\ell^1(\cos\theta), \label{eq:fg}
\end{align}
where $\delta_\ell$ and $\delta_{-\ell-1}$ are the phase shifts for
$\kappa_D=\ell$ and $\kappa_D=-\ell-1$ (sign conventions for $g$ vary between
authors). The unpolarised cross section is $|f|^2+|g|^2$. The two amplitudes
differ in phase, so the scattered beam is polarised perpendicular to the
scattering plane, $P(\theta)=2\Imag(fg^*)/(|f|^2+|g|^2)$, a spin--orbit effect
with no counterpart for spin~0 and the basis of Mott polarimetry of electron
beams. The Coulomb series \eqref{eq:fg} converge only distributionally, as in
\eqref{eq:fC-pw}, and are summed in practice by series-reduction methods
(Yennie, Ravenhall and Wilson).

For orientation, the leading term for $\zeta\ll1$, which is obtained
perturbatively and is quoted here without derivation, is Mott's correction to the
relativistic Rutherford law \eqref{eq:ruth-rel}:
\begin{equation}\label{eq:mott-factor}
  \dsdO=\Bigl(\frac{\alpha}{2pv}\Bigr)^{2}
  \frac{1-\beta^2\sin^2(\theta/2)}{\sin^4(\theta/2)},\qquad\beta=\frac vc .
\end{equation}
The extra factor tends to~1 as $\beta\to0$, and as $\beta\to1$ it becomes
$\cos^2(\theta/2)$. The next section shows that this last behaviour is exact and
does not depend on perturbation theory.

\section{Helicity and backscattering}\label{sec:helicity}

For a massless Dirac particle, or for $E\gg mc^2$, the Coulomb potential (the
time component of a vector field) couples to the right- and left-handed
components separately, so \emph{helicity is conserved}. Consider backscattering,
$\theta=\pi$, of a particle with positive helicity incident along $+z$. For a
plane wave along $z$ the orbital angular momentum component $L_z$ is zero, so
the initial state has $J_z=+\tfrac12$. In the final state the momentum is along
$-z$ and, by helicity conservation, the spin points along $-z$, so
$J_z=-\tfrac12$, again with $L_z=0$ about the axis. A central potential conserves
$J_z$, so the amplitude must vanish:
\begin{equation}\label{eq:backward}
  \dsdO\Big|_{\theta=\pi}=0\quad(\beta=1),\qquad
  \dsdO\propto\cos^2\frac\theta2\quad\text{near }\theta=\pi .
\end{equation}
This is an exact statement, reproduced by the Mott factor
\eqref{eq:mott-factor} at lowest order in $\zeta$. For $\beta<1$ helicity is not
conserved and the backward cross section is non-zero. Compare the spinless
non-relativistic law, which has $\dd\sigma/\dd\Omega(\pi)=(\alpha/4E)^2$, finite
and equal to the law's minimum.

\section{The Klein step}

A simple exactly solvable relativistic scattering problem, non-perturbative in
the strength of the potential and with a potential that does not vanish at
infinity, is the step $V(x)=0$ for $x<0$ and $V(x)=V_0$ for $x>0$, with a
particle incident from the left. Let $E>mc^2$ be the energy,
$\hbar k=\sqrt{E^2-m^2c^4}/c$ the momentum outside, and $\hbar q$ the momentum
inside, with $\hbar^2c^2q^2=(E-V_0)^2-m^2c^4$. Three regimes occur.
\begin{enumerate}
  \item $V_0<E-mc^2$. $q$ is real and the usual partial reflection and
        transmission occur.
  \item $E-mc^2<V_0<E+mc^2$. $q$ is imaginary, the wave is evanescent, and
        $R=1$.
  \item $V_0>E+mc^2$. $q$ is real again, but $E-V_0<-mc^2$: the transmitted wave
        is a negative-energy solution whose group velocity points opposite to its
        momentum. The outgoing transmitted wave has $q<0$.
\end{enumerate}
Matching the Dirac spinor at $x=0$ gives, in regimes 1 and 3,
\begin{equation}\label{eq:klein}
  R=\Bigl(\frac{1-r_K}{1+r_K}\Bigr)^{2},\quad
  T=\frac{4r_K}{(1+r_K)^2},\quad
  r_K=\frac{|q|}{k}\,\frac{E+mc^2}{|E-V_0+mc^2|},
\end{equation}
with $R+T=1$. In regime 3 both $R$ and $T$ are positive and $T$ does not vanish as
$V_0\to\infty$: it tends to $4r_K/(1+r_K)^2$ with
$r_K\to\sqrt{(E+mc^2)/(E-mc^2)}$. For $E=2mc^2$ this gives $T\approx0.93$, and
$T\to1$ as $E\gg mc^2$. A particle sees an almost transparent barrier, however
high the step. This is the Klein paradox (1929).

The resolution lies beyond single-particle quantum mechanics. A step sharp on the
Compton scale $\hbar/mc$, with $V_0>2mc^2$, can pull particle--antiparticle pairs
out of the vacuum, and the transmitted current is carried by the created
antiparticles. For a step smoothed over a distance much longer than $\hbar/mc$ the
transmission is exponentially small (Sauter, 1931); in a uniform field
$\mathcal E$ the pair-creation probability scales as
$\exp(-\pi m^2c^3/e\mathcal E\hbar)$ (Schwinger). For spin~0 the Klein regime
shows the bosonic version of the effect: the reflected flux exceeds the incident
flux, because the pair creation is stimulated.

\section{Problems}
\begin{enumerate}
  \item Show that \eqref{eq:kg} reduces to \eqref{eq:kg-rad} and
        \eqref{eq:Lambda}, and find the condition on $\ell$ and $\zeta$ for which
        the resulting radial problem is well behaved at the origin.
  \item Verify that the semiclassical rule $\Theta=2\,\dd\delta_\ell/\dd\ell$
        applied to \eqref{eq:delta-kg} gives exactly \eqref{eq:dphi-rel}.
  \item Derive \eqref{eq:kg-levels} from the poles of \eqref{eq:delta-kg} and
        show that for $\zeta\ll1$ it gives
        $E\approx mc^2[1-\zeta^2/2n^2+\dots]$, the Bohr levels with relativistic
        corrections.
  \item Derive \eqref{eq:klein} by matching the Dirac spinor at $x=0$, and show
        that the current is conserved in each regime.
  \item Explain why the helicity argument of Section~\ref{sec:helicity} fails
        for a scalar potential, which couples to the mass term $\beta mc^2$ and
        therefore flips chirality.
\end{enumerate}

\section*{Notes and references}
The Dirac equation in a Coulomb field was solved by Darwin and Gordon \cite{Darwin28,GordonA}; see also \cite{BetheSalpeter,Rose,BLP}. Mott's phase-shift treatment of the scattering of electrons is \cite{Mott29}; the first-order correction in $Z\alpha$ is due to McKinley and Feshbach \cite{McKinley48}; the reduction of the slowly convergent Coulomb series to a rapidly convergent one was given by Yennie, Ravenhall and Wilson \cite{YRW54}; and the review \cite{Motz64} collects the results and their range of validity. The polarization (Sherman) function and its use are described in \cite{Sherman56,GayDunning}. The Klein paradox is analysed in \cite{Klein29,Sauter31,Hansen81}, with a historical review in \cite{Dombey99}; Schwinger's calculation is \cite{Schwinger51}. For $Z>137$ see \cite{Pomeranchuk45,ZP71,Reinhardt77,Greiner}.

% =====================================================================
\chapter{The Invariant Formulation}
% =====================================================================

\section{The S-matrix and the invariant amplitude}

At relativistic energies particle number is not conserved, so the natural
objects are asymptotic multiparticle states and the S-matrix between them. In
units $\hbar=c=1$ and with the normalisation
$\langle p|p'\rangle=2E(2\pi)^3\delta^3(\bm p-\bm p')$, the invariant amplitude
$\mathcal M$ is defined by
\begin{equation}\label{eq:Smat}
  S_{fi}=\delta_{fi}+\ii\,(2\pi)^4\delta^4(P_f-P_i)\,\mathcal M_{fi}.
\end{equation}
This is a definition, valid whether or not $\mathcal M$ can be expanded in the
coupling. The cross section for $2\to n$ is
\begin{equation}\label{eq:dsigma}
  \dd\sigma=\frac1F\,|\mathcal M|^2\,\dd\Phi_n,\qquad
  \dd\Phi_n=(2\pi)^4\delta^4\Bigl(P_f-p_1-p_2\Bigr)\prod_{j}
  \frac{\dd^3p_j}{(2\pi)^3\,2E_j},
\end{equation}
with the invariant flux factor $F=4p^*\sqrt s$ of Section~1.3. For elastic
$2\to2$ scattering this becomes
\begin{equation}\label{eq:dsdOm-M}
  \dsdO\Big|^*=\frac{|\mathcal M|^2}{64\pi^2s},\qquad
  \frac{\dd\sigma}{\dd t}=\frac{|\mathcal M|^2}{64\pi\,s\,p^{*2}},
\end{equation}
which are mutually consistent through \eqref{eq:dsdt}. Comparison with
\eqref{eq:xs-f} identifies $f=\mathcal M/(8\pi\sqrt s)$ up to a phase convention.
Unitarity of $S$ gives the relativistic optical theorem,
\begin{equation}\label{eq:opt-rel}
  \Imag\mathcal M(s,t=0)=2p^*\sqrt s\,\sigma_{\mathrm{tot}}=\tfrac12F\,\sigma_{\mathrm{tot}},
\end{equation}
which reduces to \eqref{eq:optical} in the non-relativistic limit.

\section{Singularities, crossing and the range of the force}

The amplitude $\mathcal M(s,t,u)$ is a function of the invariants
\eqref{eq:mandelstam}. Crossing symmetry says that one analytic function
describes all the channels related by exchanging incoming and outgoing particles:
$a+b\to c+d$ (the $s$-channel), $a+\bar c\to\bar b+d$ (the $t$-channel) and so on,
each in its own physical region. Singularities of $\mathcal M$ carry the
dynamical information:
\begin{itemize}
  \item a pole in $s$ at $s=m_R^2$ is a bound state or resonance in the
        $s$-channel;
  \item a pole in $t$ at $t=m_{\mathrm{ex}}^2$ is the exchange of a particle of
        mass $m_{\mathrm{ex}}$ between the scatterers, with
        \begin{equation}\label{eq:pole}
          \mathcal M\simeq\frac{g^2}{m_{\mathrm{ex}}^2-t}
        \end{equation}
        near the pole.
\end{itemize}
In the non-relativistic limit $t=-q^2$, with $q$ the momentum transfer, and
\eqref{eq:pole} is the momentum-space form of a Yukawa force of range
$\hbar/m_{\mathrm{ex}}c$. The position of the nearest singularity in $t$ is
therefore a model-independent measure of the range of the force.

This gives a precise sense to long range. If the exchanged particle is massive,
the nearest singularity lies at $t=m_{\mathrm{ex}}^2>0$, separated from the
physical region $t\le0$ by a finite gap. Then $\dd\sigma/\dd t$ is finite at
$t=0$ and $\sigma_{\mathrm{tot}}$ is finite. If the exchanged particle is
massless, as for the photon, the pole sits at $t=0$, on the edge of the physical
region, and
\begin{equation}\label{eq:pole-xs}
  \mathcal M\propto\frac1t\quad\Longrightarrow\quad\frac{\dd\sigma}{\dd t}\propto\frac1{t^2},
\end{equation}
which is the Rutherford law \eqref{eq:dsdt-rel}. The infinite Coulomb cross
section is the integral of $1/t^2$ up to the edge of the physical region.

The same singularity controls the partial-wave expansion. For
\[
  \mathcal M=16\pi\sum_\ell(2\ell+1)\,a_\ell(s)\,P_\ell(\cos\theta),
\]
the series converges inside an ellipse in $\cos\theta$ with foci $\pm1$ (the
Lehmann--Martin ellipse) that passes through $z_t=1+t_0/(2p^{*2})$, where $t_0$
is the nearest $t$-singularity. For a massive exchange the ellipse encloses the whole physical
segment $-1\le\cos\theta\le1$. For the photon $t_0=0$, the ellipse collapses onto
the segment, and the series has no region of ordinary convergence. This is exactly
what we found in \eqref{eq:fC-pw}: the non-convergence of the Coulomb partial-wave
series is not a computational nuisance but the partial-wave image of the
masslessness of the photon.

\section{Infrared structure and the relativistic Dollard problem}

In quantum electrodynamics the Coulomb phase of Section~\ref{sec:dollard} has a
relativistic counterpart. Amplitudes for charged particles with fixed momenta
are infrared divergent, because the long-range field couples to arbitrarily soft
photons. The divergence in the cross section cancels against the emission of
real soft photons of energy below the detector resolution $\Delta E$ (Bloch and
Nordsieck, 1937), so that the physical elastic cross section depends on
$\Delta E$. The asymptotic states of charged particles are not Fock states: they
carry coherent clouds of soft photons (Chung, 1965; Kulish and Faddeev, 1970),
which play the role of Dollard's modified free dynamics \eqref{eq:dollard}.
Radiation is negligible in the non-relativistic regime, which is why
Chapter~4 can treat Coulomb scattering as strictly elastic.

\section{Thresholds and the Sommerfeld--Sakharov factor}

The penetration factor \eqref{eq:sommerfeld} survives relativistically as the
factor by which a production cross section near threshold is modified by the
Coulomb interaction of the final particles. For an oppositely charged pair of
relative velocity $v_{\mathrm{rel}}=2\beta c$ the enhancement is
\begin{equation}
  S(\beta)=\frac{2\pi|\eta|}{1-\ee^{-2\pi|\eta|}},\qquad
  2\pi|\eta|=\frac{\pi\afs}{\beta},
\end{equation}
which makes the cross section for pair production, for example
$\mathrm e^+\mathrm e^-\to\mu^+\mu^-$ or $\pi^+\pi^-$, jump to a finite value at
threshold instead of vanishing like the phase-space factor $\beta$.

\section{Problems}
\begin{enumerate}
  \item Derive \eqref{eq:dsdOm-M} from \eqref{eq:dsigma} using the two-body phase
        space $\dd\Phi_2=p^*\,\dd\Omega/(16\pi^2\sqrt s)$.
  \item Show that \eqref{eq:opt-rel} reduces to \eqref{eq:optical} when
        $f=\mathcal M/(8\pi\sqrt s)$.
  \item For $\mathcal M=g^2/(m^2-t)$ in the non-relativistic limit, compute the
        position-space potential and verify that its range is $\hbar/mc$.
\end{enumerate}

\section*{Notes and references}
The invariant amplitude and the cross-section conventions follow \cite{Weinberg,PS,Itzykson}. The analytic properties of the S-matrix and the relation between singularities and the range of forces are the subject of \cite{Eden,MartinSpearman}; the convergence of the partial-wave expansion within an ellipse is due to Lehmann \cite{Lehmann58} and, in its sharpest form, to Martin \cite{Martin63}. Infrared divergences and the Bloch--Nordsieck theorem are in \cite{BlochNordsieck,Weinberg}; the coherent dressing of charged asymptotic states is treated in \cite{Chung65,KF70}. The Sakharov--Sommerfeld enhancement is \cite{Sakharov48,Sommerfeld31}.

% =====================================================================
\chapter{Synthesis: Three Descriptions of the Same Force}
% =====================================================================

\section{The Coulomb problem in the three settings}

Two dimensionless numbers control every result of this book: $\eta=\alpha/\hbar v$,
which measures how classical the motion is, and $\zeta=\alpha/\hbar c=\eta v/c$,
which measures how relativistic it is. The three descriptions compare as follows.

\medskip
\noindent\small
\begin{tabular}{@{}>{\raggedright\arraybackslash}p{2.3cm}>{\raggedright\arraybackslash}p{3.5cm}>{\raggedright\arraybackslash}p{3.9cm}>{\raggedright\arraybackslash}p{4.8cm}@{}}
\hline
 & \textbf{Classical} & \textbf{Quantum (non-rel.)} & \textbf{Relativistic}\\
\hline
Coulomb cross section &
  $(\alpha/4E)^2/\sin^4(\theta/2)$, \eqref{eq:ruth} &
  identical, exact for every $\eta$, \eqref{eq:ruth-q} &
  leading order $(\alpha/2pv)^2/\sin^4(\theta/2)$, \eqref{eq:ruth-rel}; Mott factor for spin~$\tfrac12$, \eqref{eq:mott-factor}; exact via \eqref{eq:delta-kg}, \eqref{eq:fg}\\[6pt]
Phase of the amplitude &
  none &
  $2\sigma_0-\eta\ln\sin^2\tfrac\theta2$, \eqref{eq:fC} &
  same structure with $\ell\to\Lambda$ or $\kappa_D\to\gamma$\\[6pt]
Total cross section &
  infinite for every infinite-range $V$ &
  infinite for $V\propto r^{-1},r^{-2}$; finite iff $V$ falls faster than $r^{-2}$ &
  infinite for Coulomb (massless exchange)\\[6pt]
Bound motion &
  closed Kepler ellipses &
  Bohr levels as poles of $f_C$, \eqref{eq:bohr} &
  precessing rosettes; levels \eqref{eq:kg-levels} (spin 0) and \eqref{eq:dirac-levels} (spin $\tfrac12$)\\[6pt]
Fall to centre &
  $1/r^2$ only, \eqref{eq:cap2} &
  $1/r^2$ only, Section~\ref{sec:qinv2} &
  Coulomb, $L<\alpha/c$, \eqref{eq:cap-rel}; KG for $\zeta>\tfrac12$; none in the Dirac case for $\zeta<1$\\[6pt]
Controlling parameter &
  none (scale $\alpha/E$) &
  $\eta\gg1$ is the classical limit &
  $\zeta$ and $\eta$\\
\hline
\end{tabular}
\normalsize

\section{Five faces of a long-range force}

The same fact about the $1/r$ potential shows up in five apparently unrelated
ways.
\begin{enumerate}
  \item The incident and outgoing waves carry a logarithmic phase,
        \eqref{eq:psiC-asym}.
  \item The standard scattering theory fails and must be replaced by Dollard's,
        \eqref{eq:dollard}, with the overall phase of the amplitude
        convention-dependent.
  \item The total cross section is infinite and the angular distribution
        diverges as $\theta^{-4}$, \eqref{eq:ruth-q}.
  \item The invariant amplitude has a pole at $t=0$, on the boundary of the
        physical region, \eqref{eq:pole-xs}.
  \item The partial-wave series has no region of ordinary convergence,
        \eqref{eq:fC-pw}.
\end{enumerate}

\section{What is and is not an accident}

The coincidence of the classical and quantum Rutherford cross sections is an
accident of the $1/r$ potential and its hidden $SO(4)$ symmetry. It does not hold
for $1/r^2$, where the deflection function is exact but the cross section is not;
it does not hold for the screened potentials of real plasmas; and relativity
destroys it in two distinct ways: through the shift $\ell\to\Lambda$ in the
orbital angular momentum, which makes the orbits precess and creates a capture
region, and through spin, which adds the Mott factor and, in the ultra-relativistic
limit, forbids backscattering by helicity conservation. What survives every
change of setting is the structural content of Part~I: the definition of the cross
section, its Lorentz invariance, unitarity, and the connection between the range
of the force and the singularities of the amplitude.

\section*{Notes and references}
For comprehensive treatments of non-relativistic collision theory see \cite{MM,Newton,Taylor,GW}. The companion volume on inelastic scattering develops the consequences of coupling to other channels. Tables and numerical methods for the Coulomb functions are in \cite{AS,DLMF,Seaton02}.

\backmatter
\renewcommand{\bibname}{References}
\addcontentsline{toc}{chapter}{References}
\begin{thebibliography}{99}
\bibitem{AS} M. Abramowitz and I. A. Stegun (eds.), \emph{Handbook of Mathematical Functions}, Dover, New York (1965). Chapters 13 and 14 treat Kummer and Coulomb wave functions.
\bibitem{Bander66} M. Bander and C. Itzykson, Rev. Mod. Phys. \textbf{38}, 330 and 346 (1966).
\bibitem{Bargmann36} V. Bargmann, Z. Phys. \textbf{99}, 576 (1936).
\bibitem{BLP} V. B. Berestetskii, E. M. Lifshitz and L. P. Pitaevskii, \emph{Quantum Electrodynamics}, 2nd ed., Pergamon, Oxford (1982).
\bibitem{BerryMount} M. V. Berry and K. E. Mount, Rep. Prog. Phys. \textbf{35}, 315 (1972).
\bibitem{BetheSalpeter} H. A. Bethe and E. E. Salpeter, \emph{Quantum Mechanics of One- and Two-Electron Atoms}, Springer, Berlin (1957).
\bibitem{BlochNordsieck} F. Bloch and A. Nordsieck, Phys. Rev. \textbf{52}, 54 (1937).
\bibitem{BK} E. Byckling and K. Kajantie, \emph{Particle Kinematics}, Wiley, London (1973).
\bibitem{Case50} K. M. Case, Phys. Rev. \textbf{80}, 797 (1950).
\bibitem{Chung65} V. Chung, Phys. Rev. \textbf{140}, B1110 (1965).
\bibitem{Darwin28} C. G. Darwin, Proc. R. Soc. A \textbf{118}, 654 (1928).
\bibitem{DLMF} \emph{NIST Digital Library of Mathematical Functions}, \texttt{https://dlmf.nist.gov}. Chapters 13 and 33 treat confluent hypergeometric and Coulomb functions.
\bibitem{Dollard64} J. D. Dollard, J. Math. Phys. \textbf{5}, 729 (1964).
\bibitem{Dombey99} N. Dombey and A. Calogeracos, Phys. Rep. \textbf{315}, 41 (1999).
\bibitem{Eden} R. J. Eden, P. V. Landshoff, D. I. Olive and J. C. Polkinghorne, \emph{The Analytic S-Matrix}, Cambridge University Press (1966).
\bibitem{Fock35} V. Fock, Z. Phys. \textbf{98}, 145 (1935).
\bibitem{FordWheeler} K. W. Ford and J. A. Wheeler, Ann. Phys. (N.Y.) \textbf{7}, 259 and 287 (1959).
\bibitem{Gamow28} G. Gamow, Z. Phys. \textbf{51}, 204 (1928).
\bibitem{GayDunning} T. J. Gay and F. B. Dunning, Rev. Sci. Instrum. \textbf{63}, 1635 (1992).
\bibitem{GW} M. L. Goldberger and K. M. Watson, \emph{Collision Theory}, Wiley, New York (1964); Dover reprint (2004).
\bibitem{Goldstein} H. Goldstein, C. Poole and J. Safko, \emph{Classical Mechanics}, 3rd ed., Addison-Wesley, San Francisco (2002).
\bibitem{GordonA} W. Gordon, Z. Phys. \textbf{48}, 11 (1928).
\bibitem{GordonB} W. Gordon, Z. Phys. \textbf{48}, 180 (1928).
\bibitem{GR} I. S. Gradshteyn and I. M. Ryzhik, \emph{Table of Integrals, Series, and Products}, 7th ed., Academic Press, Amsterdam (2007).
\bibitem{Greiner} W. Greiner, B. M\"uller and J. Rafelski, \emph{Quantum Electrodynamics of Strong Fields}, Springer, Berlin (1985).
\bibitem{Hagedorn} R. Hagedorn, \emph{Relativistic Kinematics}, Benjamin, New York (1963).
\bibitem{Hansen81} A. Hansen and F. Ravndal, Phys. Scr. \textbf{23}, 1036 (1981).
\bibitem{Itzykson} C. Itzykson and J.-B. Zuber, \emph{Quantum Field Theory}, McGraw-Hill, New York (1980).
\bibitem{Jackson} J. D. Jackson, \emph{Classical Electrodynamics}, 3rd ed., Wiley, New York (1999). Chapters 13 (collisions and energy loss), 14 (radiation by moving charges) and 15 (bremsstrahlung) are the most relevant.
\bibitem{Klein29} O. Klein, Z. Phys. \textbf{53}, 157 (1929).
\bibitem{KN29} O. Klein and Y. Nishina, Z. Phys. \textbf{52}, 853 (1929).
\bibitem{KF70} P. P. Kulish and L. D. Faddeev, Theor. Math. Phys. \textbf{4}, 745 (1970).
\bibitem{LL2} L. D. Landau and E. M. Lifshitz, \emph{The Classical Theory of Fields}, 4th ed., Pergamon, Oxford (1975).
\bibitem{LL1} L. D. Landau and E. M. Lifshitz, \emph{Mechanics}, 3rd ed., Pergamon, Oxford (1976). Chapter IV treats collisions.
\bibitem{LL3} L. D. Landau and E. M. Lifshitz, \emph{Quantum Mechanics: Non-Relativistic Theory}, 3rd ed., Pergamon, Oxford (1977). Chapters XVII and XVIII treat elastic and inelastic collisions.
\bibitem{Lehmann58} H. Lehmann, Nuovo Cimento \textbf{10}, 579 (1958).
\bibitem{LS50} B. A. Lippmann and J. Schwinger, Phys. Rev. \textbf{79}, 469 (1950).
\bibitem{Mandelstam58} S. Mandelstam, Phys. Rev. \textbf{112}, 1344 (1958).
\bibitem{Martin63} A. Martin, Phys. Rev. \textbf{129}, 1432 (1963); Nuovo Cimento A \textbf{42}, 930 (1966).
\bibitem{MartinSpearman} A. D. Martin and T. D. Spearman, \emph{Elementary Particle Theory}, North-Holland, Amsterdam (1970).
\bibitem{McKinley48} W. A. McKinley and H. Feshbach, Phys. Rev. \textbf{74}, 1759 (1948).
\bibitem{Messiah} A. Messiah, \emph{Quantum Mechanics}, vols.~1--2, North-Holland, Amsterdam (1961--1962).
\bibitem{Mott28} N. F. Mott, Proc. R. Soc. A \textbf{118}, 542 (1928).
\bibitem{Mott29} N. F. Mott, Proc. R. Soc. A \textbf{124}, 425 (1929); \textbf{135}, 429 (1932).
\bibitem{Mott30} N. F. Mott, Proc. R. Soc. A \textbf{126}, 259 (1930).
\bibitem{MM} N. F. Mott and H. S. W. Massey, \emph{The Theory of Atomic Collisions}, 3rd ed., Clarendon Press, Oxford (1965).
\bibitem{Motz64} J. W. Motz, H. Olsen and H. W. Koch, Rev. Mod. Phys. \textbf{36}, 881 (1964).
\bibitem{Newton} R. G. Newton, \emph{Scattering Theory of Waves and Particles}, 2nd ed., Springer, New York (1982); Dover reprint (2002).
\bibitem{Pauli26} W. Pauli, Z. Phys. \textbf{36}, 336 (1926).
\bibitem{PS} M. E. Peskin and D. V. Schroeder, \emph{An Introduction to Quantum Field Theory}, Addison-Wesley, Reading (1995).
\bibitem{Pomeranchuk45} I. Pomeranchuk and Ya. Smorodinsky, J. Phys. (USSR) \textbf{9}, 97 (1945).
\bibitem{ReedSimonII} M. Reed and B. Simon, \emph{Methods of Modern Mathematical Physics II: Fourier Analysis, Self-Adjointness}, Academic Press, New York (1975).
\bibitem{ReedSimonIII} M. Reed and B. Simon, \emph{Methods of Modern Mathematical Physics III: Scattering Theory}, Academic Press, New York (1979).
\bibitem{Reinhardt77} J. Reinhardt and W. Greiner, Rep. Prog. Phys. \textbf{40}, 219 (1977).
\bibitem{Rose} M. E. Rose, \emph{Relativistic Electron Theory}, Wiley, New York (1961).
\bibitem{Rutherford11} E. Rutherford, Phil. Mag. \textbf{21}, 669 (1911).
\bibitem{Sakharov48} A. D. Sakharov, Zh. Eksp. Teor. Fiz. \textbf{18}, 631 (1948).
\bibitem{Sauter31} F. Sauter, Z. Phys. \textbf{69}, 742 (1931).
\bibitem{Schwinger51} J. Schwinger, Phys. Rev. \textbf{82}, 664 (1951); \textbf{91}, 728 (1953).
\bibitem{Seaton02} M. J. Seaton, Comput. Phys. Commun. \textbf{146}, 225 (2002).
\bibitem{Sherman56} N. Sherman, Phys. Rev. \textbf{103}, 1601 (1956).
\bibitem{Sommerfeld16} A. Sommerfeld, Ann. Phys. (Leipzig) \textbf{51}, 1 (1916).
\bibitem{Sommerfeld31} A. Sommerfeld, Ann. Phys. (Leipzig) \textbf{403}, 257 (1931).
\bibitem{Taylor} J. R. Taylor, \emph{Scattering Theory: The Quantum Theory of Nonrelativistic Collisions}, Wiley, New York (1972); Dover reprint (2006).
\bibitem{Weinberg} S. Weinberg, \emph{The Quantum Theory of Fields}, vol.~1, Cambridge University Press (1995). Chapter 3 treats scattering theory and Chapter 13 infrared effects.
\bibitem{YRW54} D. R. Yennie, D. G. Ravenhall and R. N. Wilson, Phys. Rev. \textbf{95}, 500 (1954).
\bibitem{ZP71} Ya. B. Zel'dovich and V. S. Popov, Usp. Fiz. Nauk \textbf{105}, 403 (1971) [Sov. Phys. Usp. \textbf{14}, 673 (1972)].
\bibitem{PDG} F. Takahashi \emph{et al.} (Particle Data Group), \emph{Review of Particle Physics}, Int. J. Mod. Phys. A \textbf{41}, 2630011 (2026), \texttt{https://pdg.lbl.gov}. The reviews on kinematics and on the passage of particles through matter are the relevant ones.
\end{thebibliography}

\end{document}


